\documentclass[10pt,a4paper,twocolumn]{article}

\usepackage[margin=0.65in]{geometry}
\usepackage[T1]{fontenc}
\usepackage{lmodern}
\usepackage{graphicx}
\usepackage{booktabs}
\usepackage{hyperref}
\usepackage{enumitem}
\usepackage{array}
\usepackage{xcolor}
\usepackage{caption}
\usepackage{subcaption}
\usepackage{placeins}
\usepackage{fontawesome5}

\hypersetup{
    colorlinks=true,
    linkcolor=blue!50!black,
    urlcolor=blue!50!black,
    citecolor=blue!50!black
}

\setlength{\columnsep}{0.22in}
\setlength{\parindent}{0pt}
\setlength{\parskip}{0.35em}
\captionsetup{font=small}
\captionsetup[figure]{skip=4pt}
\setlength{\textfloatsep}{8pt plus 2pt minus 2pt}
\setlength{\floatsep}{8pt plus 2pt minus 2pt}
\setlength{\intextsep}{8pt plus 2pt minus 2pt}
\setcounter{topnumber}{2}
\setcounter{dbltopnumber}{2}
\renewcommand{\topfraction}{0.9}
\renewcommand{\dbltopfraction}{0.9}
\renewcommand{\textfraction}{0.07}
\renewcommand{\floatpagefraction}{0.8}
\renewcommand{\dblfloatpagefraction}{0.8}

\title{\textbf{RepViT on Low-Resolution Datasets: A CIFAR Study}}
\author{Vaibhav Garg \and Aaditya Bansal \and Rudra Gupta}
\date{March 28, 2026}

\begin{document}
\maketitle

\begin{center}
\small \href{https://github.com/rudra-dotcom/DL-Project.git}{\faGithub\ GitHub Repository}
\hspace{1.5em}
\href{<ADD-YOUTUBE-LINK-HERE>}{\faYoutube\ Demo Video}
\end{center}

\begin{abstract}
This project studies how RepViT, a mobile-oriented convolutional architecture proposed in \emph{RepViT: Revisiting Mobile CNN From ViT Perspective}, behaves when transferred from ImageNet-scale images to low-resolution CIFAR benchmarks. The central observation is that RepViT's original early downsampling is well suited to \(224\times224\) inputs but too aggressive for \(32\times32\) images, causing early spatial information loss. To address this, the project develops three complementary contributions: a low-resolution macro redesign for CIFAR-100, an attention-and-training-recipe redesign for CIFAR-100, and a CIFAR-10-specific RepViT variant with an added inference and deployment workflow. Across the three efforts, the project shows that preserving spatial resolution longer, placing channel recalibration more selectively, and pairing architectural changes with task-appropriate training recipes produce large gains over the unmodified baseline.
\end{abstract}

\section{Base/Selected Paper Overview}

RepViT is a family of lightweight \emph{pure CNN} models introduced by Wang et al. in the CVPR 2024 paper \emph{RepViT: Revisiting Mobile CNN From ViT Perspective} \cite{repvit}. The paper revisits MobileNet-style convolutional models from the perspective of efficient Vision Transformers and argues that much of the mobile performance gap often attributed to self-attention actually comes from better architecture design choices rather than attention alone.

The base paper starts from MobileNetV3 and incrementally improves three design levels:
\begin{itemize}[leftmargin=1.2em]
    \item \textbf{Block structure:} revisiting how local mixing, channel mixing, and residual connections are arranged.
    \item \textbf{Macro design:} organizing stages and downsampling in a way that better matches mobile efficiency constraints.
    \item \textbf{Micro design:} using implementation details that improve latency on deployment hardware.
\end{itemize}

The contribution of the original paper is therefore not a single isolated module, but a mobile CNN family that borrows the most effective efficiency ideas from lightweight ViTs while staying convolutional. The paper reports that RepViT achieves over \(80\%\) ImageNet top-1 accuracy with roughly \(1\) ms latency on iPhone 12, establishing a strong accuracy-latency tradeoff for mobile vision workloads.

\subsection*{RepViT Architecture}

The vanilla RepViT pipeline can be summarized at a high level as:
\begin{itemize}[leftmargin=1.2em]
    \item an initial \textbf{patch embedding stem} with aggressive early downsampling,
    \item a sequence of mobile-friendly \textbf{RepViT blocks} organized into stages,
    \item selective use of \textbf{channel recalibration} modules such as SE,
    \item progressive downsampling and channel expansion across stages,
    \item a lightweight classification head for ImageNet-scale prediction.
\end{itemize}

This vanilla architecture is the starting point for all three contributions in the project. Each contribution modifies this baseline differently to make it more compatible with CIFAR-sized inputs.

\begin{figure}[!htbp]
    \centering
    \includegraphics[width=0.88\columnwidth]{repvit_architecture.png}
    \caption{RepViT Architecture}
    \label{fig:vanilla-block}
\end{figure}


For this project, the most important background detail is that RepViT was designed and benchmarked primarily for large ImageNet inputs (\(224\times224\)). Its patch embedding stem contains consecutive stride-2 convolutions, which is reasonable at high resolution but much more destructive for CIFAR images. On \(32\times32\) inputs, the feature map shrinks to \(8\times8\) almost immediately. That mismatch motivates the entire project: if RepViT is strong on mobile hardware, can it be re-designed to behave better on low-resolution datasets without giving up efficiency?

\section{Our Approach and Contribution}

Rather than proposing one single modification, the project contributes three complementary responses to the same root problem: \textbf{RepViT is overly aggressive in its early spatial compression for CIFAR-scale inputs.}

\subsection*{1.Rudra: Low-Resolution Macro Design}

Rudra's contribution focuses on CIFAR-100 with RepViT-M0.9. The baseline model achieves only \(43.07\%\) top-1, showing that the original macro design is poorly matched to \(32\times32\) inputs. Two new variants were introduced:
\begin{itemize}[leftmargin=1.2em]
    \item \texttt{repvit\_m0\_9\_lr}: changes the stem stride from \((2,2)\) to \((1,2)\) and removes the final downsampling step.
    \item \texttt{repvit\_m0\_9\_lr\_rase}: keeps the low-resolution macro design and adds a \textbf{resolution-aware SE schedule}, assigning more SE emphasis to early high-resolution stages and less to later low-resolution stages.
\end{itemize}

This contribution is important because it directly attacks the spatial-collapse problem at the architectural level. The design hypothesis is simple and strong: for small images, preserve resolution longer and spend recalibration budget where spatial detail is still present.

\subsection*{2.Vaibhav: Selective Attention + Stronger Training Recipe}

Vaibhav's contribution also targets CIFAR-100 with RepViT-M0.9 but takes a different path. Instead of only changing the downsampling schedule, it redesigns channel attention and the training recipe:
\begin{itemize}[leftmargin=1.2em]
    \item the first stem convolution is changed to stride \(1\),
    \item SE blocks are replaced with \textbf{ECA} attention,
    \item attention placement becomes stage-aware: dense in early layers, sparse in middle layers, removed in the final stage,
    \item a superclass auxiliary head is added during training,
    \item the optimization pipeline is strengthened using RandAugment, MixUp, CutMix, warmup, and gradient clipping.
\end{itemize}

This contribution shows that the low-resolution issue is not only about spatial resolution. It also involves \emph{where} and \emph{how} attention-like recalibration is applied, and how training-time regularization can help small-image transfer.

\subsection*{3.Aaditya: CIFAR-10 Variant + Inference Workflow}

Aaditya's contribution shifts to CIFAR-10 using RepViT-M1.1 and builds a dedicated \texttt{repvit\_cifar10} variant. The changes are targeted and latency-aware:
\begin{itemize}[leftmargin=1.2em]
    \item patch embedding stride is reduced from \(2\) to \(1\),
    \item SE blocks are retained in alternating fashion in Stages 1 and 2,
    \item SE blocks are removed entirely from Stages 3 and 4.
\end{itemize}

The CIFAR-10 workspace also includes benchmarking, automated visual generation, and a practical inference/deployment pipeline:
\begin{itemize}[leftmargin=1.2em]
    \item single-image inference,
    \item latency and throughput measurement,
    \item Core ML export,
    \item an iOS SwiftUI demo app for on-device classification.
\end{itemize}

Taken together, the full project contribution is broader than a model tweak. It covers \textbf{architecture adaptation, training strategy, benchmarking, visualization, and deployment}, all centered on making RepViT better suited for low-resolution data.

\section{Experiment Results and Visual Outputs}

\subsection*{Quantitative Summary}

\begin{table}[h]
\centering
\caption{Best reported results from the three project contributions.}
\label{tab:summary}
\begin{tabular}{>{\raggedright\arraybackslash}p{2.0cm} >{\raggedright\arraybackslash}p{1.3cm} >{\raggedright\arraybackslash}p{2.5cm} c}
\toprule
Contributor & Dataset & Best Variant & Top-1 \\
\midrule
Rudra & CIFAR-100 & RepViT-M0.9-LR-RASE & 56.22\% \\
Vaibhav & CIFAR-100 & RepViT-M0.9 + ECA + recipe & 74.71\% \\
Aaditya & CIFAR-10 & RepViT-CIFAR10 & 91.02\% \\
\bottomrule
\end{tabular}
\end{table}

The three contributions support a common conclusion: \textbf{RepViT benefits significantly from low-resolution-aware redesign.} However, they differ in how far they push the idea:
\begin{itemize}[leftmargin=1.2em]
    \item Rudra shows that macro-level changes alone already improve CIFAR-100 performance substantially, from \(43.07\%\) to \(56.22\%\), while even reducing parameter count slightly in the best variant.
    \item Vaibhav demonstrates the strongest CIFAR-100 result at \(74.71\%\), suggesting that selective attention plus a richer training recipe is more effective than macro changes alone.
    \item Aaditya obtains the best overall accuracy in the repository on CIFAR-10, improving the RepViT-M1.1 baseline from \(87.39\%\) to \(91.02\%\) with nearly unchanged latency and fewer parameters.
\end{itemize}

\subsection*{Per-Contribution Findings}

\textbf{Rudra.} The best low-resolution-aware model (\texttt{repvit\_m0\_9\_lr\_rase}) outperforms the baseline by \(+13.15\) points on CIFAR-100. This validates the core hypothesis that preserving spatial detail longer is essential for small-image transfer.

\textbf{Vaibhav.} The ECA-based variant reaches \(74.71\%\) top-1 and \(92.18\%\) top-5 on CIFAR-100 with roughly \(3.7\) ms latency. This is the strongest CIFAR-100 result in the project and shows the value of jointly optimizing architecture and training recipe.

\textbf{Aaditya.} The CIFAR-10-specific variant improves best test accuracy by \(+3.63\) points over RepViT-M1.1 while reducing parameters from \(7.78\)M to \(4.52\)M. Its latency remains essentially unchanged (\(5.98\) ms vs. \(5.95\) ms), which is a particularly strong efficiency result.

\subsection*{Rudra's Visual Outputs}

\begin{figure}[!htbp]
    \centering
    \fbox{\parbox[c][1.45in][c]{0.88\columnwidth}{\centering Placeholder for Rudra Visual 1\\[0.4em]
    Suggested content: CIFAR-100 accuracy comparison for baseline, LR, and LR-RASE variants.}}
    \caption{Primary CIFAR-100 comparison figure placeholder.}
    \label{fig:rudra-1}
\end{figure}

\begin{figure}[!htbp]
    \centering
    \includegraphics[width=0.88\columnwidth]{rudra.png.jpeg}
    \caption{Latency visualization from Rudra's benchmark workflow.}
    \label{fig:rudra-2}
\end{figure}

\begin{figure}[!htbp]
    \centering
    \fbox{\parbox[c][1.45in][c]{0.88\columnwidth}{\centering Placeholder for Rudra Visual 3\\[0.4em]
    Suggested content: architecture diagram, feature-map comparison, or benchmark summary table for Rudra's CIFAR-100 variants.}}
    \caption{Rudra architecture or benchmark summary placeholder.}
    \label{fig:rudra-3}
\end{figure}

Rudra's visuals should emphasize that low-resolution-aware macro changes substantially improve CIFAR-100 performance and that the best variant benefits from both delayed downsampling and resolution-aware SE placement.

\vspace{4\baselineskip}

\FloatBarrier
\subsection*{Vaibhav's Visual Outputs}

\begin{figure}[!htbp]
    \centering
    \includegraphics[width=0.97\columnwidth]{vaibhav_architecture_diagram.png}
    \caption{ Vaibhav's architecture visualization.}
    \label{fig:vaibhav-3}
\end{figure}

\begin{figure}[!htbp]
    \centering
    \includegraphics[width=0.88\columnwidth]{val_acc1_vs_acc5.png}
    \caption{Validation accuracy curves.}
    \label{fig:vaibhav-1}
\end{figure}

\begin{figure}[!htbp]
    \centering
    \includegraphics[width=0.88\columnwidth]{latency.png}
    \caption{Latency plot.}
    \label{fig:vaibhav-2}
\end{figure}
\FloatBarrier

Vaibhav's outputs should communicate that stronger CIFAR-100 performance comes from the combination of selective ECA placement and a more effective training recipe, not from any single isolated change.

\FloatBarrier
\subsection*{Aaditya's Visual Outputs}

\begin{figure}[!htbp]
    \centering
    \includegraphics[width=0.88\columnwidth]{aadi_architecture.png}
    \caption{CIFAR-10 architecture or comparison visualization.}
    \label{fig:aaditya-1}
\end{figure}

\begin{figure}[!htbp]
    \centering
    \includegraphics[width=0.88\columnwidth]{aadi_result.png}
    \caption{Accuracy-latency tradeoff for Aaditya's CIFAR-10 study.}
    \label{fig:aaditya-2}
\end{figure}

\begin{figure}[!htbp]
    \centering
    \includegraphics[width=0.97\columnwidth]{aadi_summary.png}
    \caption{Summary}
    \label{aadi-3}
\end{figure}

\FloatBarrier
\subsection*{Inference App UI Placeholders}

\begin{figure}[!htbp]
    \centering
    \includegraphics[width=0.88\columnwidth]{ui_1.jpeg}
    \caption{Inference app UI: image selection or home screen.}
    \label{fig:aadi-4}
\end{figure}

\begin{figure}[!htbp]
    \centering
    \includegraphics[width=0.88\columnwidth]{ui_2.PNG}
    \caption{Inference app UI: on-device prediction result screen.}
    \label{fig:aadi-5}
\end{figure}
\FloatBarrier

Aaditya's figures show that the CIFAR-specific RepViT variant clearly outperforms the unmodified baseline while keeping latency almost unchanged, making it the strongest efficiency-preserving result in the project.

\subsection*{Overall Interpretation}

Across all experiments, three patterns repeat:
\begin{itemize}[leftmargin=1.2em]
    \item \textbf{Preserve early spatial resolution.} Every successful variant weakens or delays aggressive downsampling.
    \item \textbf{Use attention selectively.} Whether through SE scheduling or ECA replacement, attention works best when concentrated in stages where feature maps still contain useful detail.
    \item \textbf{Pair architecture with recipe.} Vaibhav's contribution shows that low-resolution adaptation is strongest when model design and training strategy are improved together.
\end{itemize}

The project therefore delivers a clear engineering conclusion: RepViT is a strong mobile vision backbone, but its default ImageNet-oriented design must be adjusted when moving to low-resolution datasets such as CIFAR-10 and CIFAR-100.

\begin{thebibliography}{4}
\bibitem{repvit}
Ao Wang, Hui Chen, Zijia Lin, Jungong Han, and Guiguang Ding.
\newblock RepViT: Revisiting Mobile CNN From ViT Perspective.
\newblock In \emph{Proceedings of the IEEE/CVF Conference on Computer Vision and Pattern Recognition}, 2024.
\newblock \url{https://arxiv.org/abs/2307.09283}

\bibitem{mobilenetv3}
Andrew Howard, Mark Sandler, Grace Chu, Liang-Chieh Chen, Bo Chen, Mingxing Tan, Weijun Wang, Yukun Zhu, Ruoming Pang, Vijay Vasudevan, Quoc V. Le, and Hartwig Adam.
\newblock Searching for MobileNetV3.
\newblock \emph{CoRR}, abs/1905.02244, 2019.
\newblock \url{https://arxiv.org/abs/1905.02244}

\bibitem{eca}
Qilong Wang, Banggu Wu, Pengfei Zhu, Peihua Li, Wangmeng Zuo, and Qinghua Hu.
\newblock ECA-Net: Efficient Channel Attention for Deep Convolutional Neural Networks.
\newblock \emph{CoRR}, abs/1910.03151, 2019.
\newblock \url{https://arxiv.org/abs/1910.03151}

\bibitem{cifar}
Alex Krizhevsky.
\newblock Learning Multiple Layers of Features from Tiny Images.
\newblock Technical Report, University of Toronto, 2009.
\newblock \url{https://www.cs.toronto.edu/~kriz/learning-features-2009-TR.pdf}
\end{thebibliography}

\end{document}
